\section{DeSFAM: Lọc Syscall Động Dựa trên eBPF và Giảm Thiểu Bất Thường}
\label{sec:desfam}

Phần này trình bày thiết kế và triển khai của framework DeSFAM được đề xuất, bao gồm
kiến trúc, các giai đoạn hoạt động và chiến lược thực thi.

\subsection{Tổng Quan}
\label{subsec:overview}

DeSFAM (Dynamic eBPF-Driven Syscall Filtering and Anomaly Mitigation) là một framework nhẹ,
hoạt động thời gian thực, được thiết kế để thu hẹp bề mặt tấn công kernel và phát hiện hành vi
syscall bất thường trong các môi trường container. Framework thực thi nguyên tắc đặc quyền tối
thiểu trong sử dụng syscall nhằm giảm thiểu các mối đe dọa như leo thang đặc quyền, thoát
container và chèn syscall.
Hình~\ref{fig:architecture} minh họa quy trình kiến trúc của framework.

DeSFAM hoạt động theo ba giai đoạn tích hợp:
\begin{itemize}
  \item \textbf{Lập Danh Sách Truy Cập Syscall Kết Hợp và Tạo Hồ Sơ Seccomp:} Kết hợp trích xuất
        đồ thị lời gọi tĩnh và theo dõi eBPF động để xây dựng danh sách truy cập syscall được
        tinh chỉnh, loại bỏ các syscall không cần thiết và liên quan đến CVE.
  \item \textbf{Giám Sát Thời Gian Thực và Phát Hiện Bất Thường (SyscallAD):} Theo dõi các
        chuỗi syscall thời gian thực và phân tích chúng bằng tập hợp VAE + Isolation Forest để
        phát hiện các bất thường ở mức chuỗi, bao gồm cả việc sử dụng độc hại các syscall được
        phép.
  \item \textbf{Thực Thi Thích Ứng và Quản Lý Chính Sách:} Thực thi việc chặn, giới hạn tốc độ
        hoặc ghi nhật ký các syscall có rủi ro cao bằng eBPF maps và LSM hooks, dựa trên điểm
        rủi ro từ SyscallAD, mà không làm gián đoạn quá trình thực thi container.
\end{itemize}

Bằng cách tích hợp lập hồ sơ nhận biết chuỗi, phát hiện bất thường và thực thi thời gian thực,
DeSFAM cung cấp khả năng bảo vệ có thể mở rộng và thích ứng, phù hợp với các triển khai
container cloud-native hiện đại.

\begin{figure}[htbp]
  \centering
  \includegraphics[width=\linewidth]{figure1}
  \caption{Kiến trúc DeSFAM thể hiện sự tương tác giữa lập hồ sơ syscall ở không gian người dùng,
           phát hiện bất thường và thực thi ở không gian kernel sử dụng eBPF và seccomp cho
           điều khiển syscall thích ứng, thời gian thực.}
  \label{fig:architecture}
\end{figure}

\subsection{Các Cân Nhắc Thiết Kế}
\label{subsec:design}

Thiết kế của DeSFAM được định hướng bởi nhu cầu thực thi nguyên tắc đặc quyền tối thiểu cho
syscall trong khi giải quyết các hạn chế của lọc chỉ tĩnh, theo dõi xuất xứ chi phí cao và
các bộ phát hiện bất thường phụ thuộc vào việc huấn luyện lại. Các hồ sơ seccomp truyền thống
hoặc quá dễ dãi hoặc quá cứng nhắc, không thể thích nghi với hành vi thời gian thực đang thay
đổi. DeSFAM khắc phục điều này bằng cách kết hợp phân tích tĩnh và động, cho phép lập hồ sơ
syscall chính xác và nhận thức ngữ cảnh.

Để đảm bảo thực thi thời gian thực mà không ảnh hưởng đến tính khả dụng của container, DeSFAM
tận dụng các eBPF probes và LSM hooks, hoạt động hoàn toàn trong không gian kernel và hỗ trợ
cập nhật chính sách động với độ trễ tối thiểu. Lựa chọn thiết kế này tránh chi phí hiệu năng
của việc xây dựng đồ thị xuất xứ liên tục và loại bỏ nhu cầu sửa đổi kernel hoặc khởi động lại
container.

Đối với phát hiện bất thường, DeSFAM áp dụng mô hình kết hợp (VAE + Isolation Forest) nắm bắt
các mẫu hành vi trong chuỗi syscall trong khi tránh huấn luyện lại thường xuyên, hỗ trợ phát
hiện có thể mở rộng và mạnh mẽ các cuộc tấn công zero-day. Không giống các giải pháp bị ràng
buộc phần cứng như TEEs, DeSFAM vẫn hoàn toàn dựa trên phần mềm và có thể di chuyển, đảm bảo
tính tương thích trên nhiều triển khai cloud-native khác nhau.

\subsection{Giai Đoạn 1: Lập Danh Sách Truy Cập Syscall Kết Hợp và Tạo Hồ Sơ Seccomp}
\label{subsec:phase1}

Bảo mật container phụ thuộc vào kiểm soát syscall chính xác, vì các container dùng chung kernel
của máy chủ và đặt ra rủi ro truy cập kernel trái phép. Giai đoạn~1 của DeSFAM giảm thiểu điều
này bằng cách xây dựng danh sách truy cập syscall dành riêng cho container thông qua phương pháp
kết hợp kết hợp phân tích tĩnh, lập hồ sơ động và thông tin tình báo mối đe dọa dựa trên CVE để
tạo ra hồ sơ seccomp tối thiểu và chính xác.

\begin{figure}[htbp]
  \centering
  \includegraphics[width=\linewidth]{figure2}
  \caption{DeSFAM Giai Đoạn~1: Lập danh sách truy cập syscall kết hợp và tạo hồ sơ seccomp
           sử dụng phân tích tĩnh, giám sát eBPF và lọc thông tin CVE để tạo ra các chính sách
           syscall an toàn dành riêng cho container.}
  \label{fig:phase1}
\end{figure}

\subsubsection{Phân Tích Tĩnh}
\label{subsubsec:static_analysis}

Phân tích tĩnh bắt đầu bằng cách kiểm tra các image container, bao gồm Dockerfile, binary và
các thư viện được liên kết, mà không thực thi chúng. Sử dụng các công cụ như \texttt{objdump}
và \texttt{readelf}, các binary ELF được dịch ngược để trích xuất các lệnh syscall cấp thấp và
các hàm wrapper của libc.

Để tăng cường độ bao phủ, một Đồ Thị Lời Gọi Tĩnh (SCG) được xây dựng bằng NetworkX để theo dõi
ánh xạ từ hàm sang syscall trên toàn cơ sở mã. Tập syscall tĩnh kết quả,
$S_{\text{static}}$, nắm bắt tất cả các syscall có thể xác định được từ luồng điều khiển của mã.

\begin{itemize}
  \item \textbf{Đầu vào:} Image container (Dockerfile, ELF binary)
  \item \textbf{Công cụ:} \texttt{objdump}, \texttt{readelf}, \texttt{networkx}
  \item \textbf{Quá trình:} Dịch ngược ELF binary, trích xuất đường dẫn syscall qua duyệt SCG
  \item \textbf{Đầu ra:} Danh sách syscall tĩnh ($S_{\text{static}}$)
\end{itemize}

Tuy nhiên, phân tích tĩnh không thể phát hiện các syscall được gọi bởi các thư viện được giải
quyết lúc chạy hoặc mã được biên dịch JIT. Do đó, trong khi $S_{\text{static}}$ cung cấp một
tập syscall nền tảng, nó phải được tinh chỉnh thông qua phân tích động trong các giai đoạn tiếp
theo.

\subsubsection{Lập Hồ Sơ Động: Nắm Bắt Hành Vi Thời Gian Thực}
\label{subsubsec:dynamic_profiling}

Để giải quyết các hạn chế của phân tích tĩnh, DeSFAM kết hợp lập hồ sơ động bằng cách giám sát
hoạt động syscall trực tiếp trong quá trình thực thi container dưới các khối lượng công việc đại
diện chuẩn hóa. Giai đoạn này nắm bắt các syscall được kích hoạt bởi các điều kiện thời gian
thực bao gồm các thư viện được tải động, logic đặc thù môi trường và các tương tác phụ thuộc,
đồng thời đảm bảo độ bao phủ toàn diện các hành vi vận hành điển hình.

Các container thực thi dưới các khối lượng công việc dành riêng cho vai trò được thiết kế bằng
các benchmark tiêu chuẩn ngành phản ánh các hồ sơ vận hành thực tế:

\begin{itemize}
  \item \textbf{Máy chủ web (Nginx):} \texttt{wrk} với 500--10,000 kết nối đồng thời;
        ApacheBench mô phỏng các yêu cầu GET/POST tại 1K--10K req/sec.
  \item \textbf{Cơ sở dữ liệu (MySQL/Redis/MongoDB):} Khối lượng công việc SysBench OLTP
        (100--500 luồng, 1M+ giao dịch); \texttt{redis-benchmark} tại 1M ops/sec.
  \item \textbf{Thời gian chạy ứng dụng (Node.js/Python):} Kiểm tra tải Locust mô phỏng 50--500
        luồng người dùng đồng thời; xử lý tập dữ liệu thực tế (phân tích JSON/CSV quy mô GB).
\end{itemize}

Các khối lượng công việc này bao gồm các mẫu lưu lượng sản xuất ở phân vị thứ 90, chuỗi khởi
động/tắt container, tương tác phụ thuộc và các kịch bản lỗi.

Trong quá trình thực thi, các eBPF probes được gắn vào các điểm vào syscall
(\texttt{sys\_enter\_*}) để theo dõi các sự kiện syscall. Siêu dữ liệu được thu thập bao gồm ID
syscall, ID tiến trình, dấu thời gian và định danh namespace.

\begin{itemize}
  \item \textbf{Đầu vào:} Các instance container dưới các khối lượng công việc chuẩn hóa
  \item \textbf{Công cụ:} eBPF qua \texttt{bcc}/\texttt{bpftrace} với theo dõi ở không gian kernel
  \item \textbf{Quá trình:} Triển khai container $\to$ gắn eBPF probes $\to$ ghi siêu dữ liệu
        vào ring buffers $\to$ tổng hợp syscall qua 3 chu kỳ khối lượng công việc
  \item \textbf{Đầu ra:} Danh sách syscall động ($S_{\text{dynamic}}$)
\end{itemize}

\subsubsection{Thông Tin Tình Báo Mối Đe Dọa: Lọc Syscall Dựa trên CVE}
\label{subsubsec:cve_filtering}

Để tăng cường nhận thức bảo mật, DeSFAM tích hợp thông tin tình báo mối đe dọa bằng cách kết
hợp dữ liệu CVE từ Cơ Sở Dữ Liệu Lỗ Hổng Quốc Gia (NVD). Sử dụng xử lý ngôn ngữ tự nhiên
(NLP), chúng tôi trích xuất tên syscall (ví dụ: \texttt{ptrace}, \texttt{clone},
\texttt{execve}) từ các mô tả CVE để xây dựng danh sách syscall có rủi ro cao, ký hiệu là
$S_{\text{blocked}}$.

Để xây dựng $S_{\text{blocked}}$, chúng tôi sử dụng một công cụ phân tích dựa trên NLP phân
tích các báo cáo CVE được lấy từ NVD qua REST API của nó. Kỹ thuật trích xuất từ khóa dựa trên
quy tắc được áp dụng để xác định tên syscall được đề cập trong các mô tả CVE. Chúng tôi sử dụng
từ điển được chọn lọc của tên syscall Linux có nguồn gốc từ các header nguồn kernel (ví dụ:
\texttt{unistd64.h}) để đảm bảo ánh xạ chính xác. Hơn nữa, để tránh chặn quá mức, các siêu dữ
liệu ngữ cảnh như khả năng áp dụng phiên bản kernel, các hệ thống con bị ảnh hưởng, điều kiện
khai thác và ràng buộc tham số được trích xuất từ văn bản CVE bằng cách khớp mẫu và phân tích
biểu thức chính quy.

\begin{itemize}
  \item \textbf{Đầu vào:} Các mục CVE có đề cập syscall
  \item \textbf{Ánh xạ:} Trích xuất syscall dựa trên quy tắc + lọc ngữ cảnh
  \item \textbf{Đầu ra:} Danh sách syscall bị chặn ($S_{\text{blocked}}$)
\end{itemize}

Tập này được loại trừ khỏi hợp của các tập syscall tĩnh và động, đảm bảo các syscall dễ bị tổn
thương bị bỏ qua ngay cả khi được quan sát trong quá trình lập hồ sơ.

\subsubsection{Tổng Hợp Danh Sách Truy Cập Cuối Cùng và Tạo Chính Sách}
\label{subsubsec:synthesis}

Danh sách truy cập cuối cùng $S_{\text{final}}$ được tính như sau:
\begin{equation}
  S_{\text{final}} = \bigl(S_{\text{static}} \cup S_{\text{dynamic}}\bigr)
                      \cap S_{\text{template}}
                      \setminus S_{\text{blocked}}
                      \setminus S_{\text{template}}
  \label{eq:sfinal}
\end{equation}

Ở đây, $S_{\text{template}}$ định nghĩa hành vi syscall mong đợi dành riêng cho vai trò của
container (ví dụ: máy chủ web, cơ sở dữ liệu, khối lượng công việc GPU), đóng vai trò là ranh
giới chức năng. Các syscall template được suy ra bằng cách kết hợp nhiều nguồn: hồ sơ seccomp
mặc định chính thức của Docker, tài liệu syscall kernel Linux, các quy tắc bảo mật do cộng đồng
duy trì từ dự án Falco, và các dấu vết syscall thực nghiệm được thu thập trong quá trình thực
thi container lành mạnh ban đầu.

Phương pháp đa nguồn này đảm bảo rằng các syscall hiếm khi được gọi nhưng cần thiết trong quá
trình khởi tạo container hoặc các đường dẫn thực thi không phổ biến được đưa vào $S_{\text{template}}$
trong khi vẫn duy trì mức tiếp xúc syscall tối thiểu.

Tập syscall dựa trên CVE $S_{\text{blocked}}$ được tạo ra bằng cách trích xuất các điều kiện
khai thác, bao gồm phiên bản kernel, các lỗ hổng dành riêng cho thiết bị, hoặc các mẫu tham số
syscall cụ thể. Điều này ngăn chặn việc chặn quá mức các syscall chỉ dễ bị tổn thương trong các
cấu hình cụ thể không áp dụng cho môi trường thời gian chạy của container.

Danh sách truy cập kết quả $S_{\text{final}}$ được biên dịch thành chính sách seccomp
(\texttt{seccomp\_profile.json}) được thực thi bởi kernel Linux lúc chạy. Hình~\ref{fig:phase1}
và Thuật toán~\ref{alg:phase1} tóm tắt quá trình này.

\begin{algorithm}[htbp]
\caption{Lập Hồ Sơ Syscall Kết Hợp Nâng Cao}
\label{alg:phase1}
\begin{algorithmic}[1]
  \Require Image container, instance đang chạy, cơ sở dữ liệu CVE
  \Ensure Danh sách truy cập syscall an toàn và có chức năng
  \State Trích xuất syscall qua phân tích tĩnh $\to S_{\text{static}}$
  \State Lập hồ sơ thời gian thực dưới khối lượng công việc (lập hồ sơ động) $\to S_{\text{dynamic}}$
  \State Phân tích CVE với phân tích ngữ cảnh $\to S_{\text{blocked}}$
  \State Tạo $S_{\text{final}} = (S_{\text{static}} \cup S_{\text{dynamic}})
         \cap S_{\text{template}} \setminus S_{\text{blocked}} \setminus S_{\text{template}}$
  \State \Return $S_{\text{final}}$
\end{algorithmic}
\end{algorithm}

\subsubsection{Tải Chính Sách Dành Riêng cho Container}
\label{subsubsec:policy_upload}

DeSFAM tải động các chính sách syscall dành riêng cho container bằng \texttt{policy\_loader}.
Mỗi chính sách phản ánh hồ sơ hành vi của container từ phân tích trước đó và được ghi vào
eBPF map được ghim tại \texttt{/sys/fs/bpf/blocked\_syscalls}. Định danh namespace đóng vai trò
là khóa để cách ly container.

Thiết kế này thực thi truy cập syscall đặc quyền tối thiểu cho mỗi container lúc chạy, mà không
cần khởi động lại instance. Hình~\ref{fig:policy_upload} cho thấy quy trình tải và việc điền vào
Syscall map.

\begin{figure}[htbp]
  \centering
  \includegraphics[width=\linewidth]{figure3}
  \caption{DeSFAM Giai Đoạn~3: (a) Tải động các chính sách syscall dành riêng cho container
           sử dụng \texttt{policy\_loader}; (b) Xác minh các syscall bị chặn trong eBPF map
           được ghim sử dụng \texttt{bpftool}.}
  \label{fig:policy_upload}
\end{figure}

\subsection{Giai Đoạn 2: Giám Sát Thời Gian Thực và Phát Hiện Bất Thường Syscall (SyscallAD)}
\label{subsec:phase2}

Giai đoạn~2 của DeSFAM giới thiệu SyscallAD, một module phát hiện bất thường thời gian thực kết
hợp học sâu và các phương pháp tập hợp để xác định cả các mối đe dọa đã biết và zero-day trong
các môi trường container. Nó giám sát các chuỗi syscall bằng eBPF và phát hiện các độ lệch hành
vi liên quan đến leo thang đặc quyền, thoát container hoặc chèn syscall.

SyscallAD sử dụng một tập hợp kết hợp gồm Bộ Mã Hóa Biến Phân (VAE), học các biểu diễn ẩn của
các mẫu syscall lành mạnh, và Isolation Forest (iForest), xác định các điểm ngoại lệ thống kê.
Sự kết hợp này cho phép phát hiện mạnh mẽ, độ trễ thấp với ít dương tính giả và chi phí tối
thiểu.

\begin{figure}[htbp]
  \centering
  \includegraphics[width=\linewidth]{figure4}
  \caption{DeSFAM Giai Đoạn~2: Giám sát syscall thời gian thực và phát hiện bất thường sử dụng
           mô hình kết hợp VAE + iForest. Module ngoại tuyến được huấn luyện trên các chuỗi
           DongTing, trong khi module trực tuyến xử lý các syscall trực tiếp qua eBPF để đánh
           giá bất thường và leo thang chính sách.}
  \label{fig:phase2}
\end{figure}

\subsubsection{Bộ Dữ Liệu DongTing: Đường Cơ Sở Hành Vi cho Phát Hiện}
\label{subsubsec:dongting}

SyscallAD được huấn luyện trên bộ dữ liệu DongTing~\cite{ref26_dongtingdataset}, gồm
18,966 chuỗi syscall và hơn 2,6~triệu syscall riêng lẻ được thu thập từ các ứng dụng container.
Bộ dữ liệu được chia thành:

\begin{itemize}
  \item \textbf{6,850 chuỗi lành mạnh:} Được thu thập từ các khối lượng công việc bình thường
        như máy chủ web và cơ sở dữ liệu thực hiện I/O tệp (\texttt{open}, \texttt{read}) và
        các thao tác bộ nhớ (\texttt{mmap}, \texttt{futex}).
  \item \textbf{12,116 chuỗi độc hại:} Được tạo ra thông qua các cuộc tấn công mô phỏng, bao
        gồm leo thang đặc quyền (\texttt{setuid}), giả mạo hệ thống tệp (\texttt{unlink}) và
        xâm nhập mạng (\texttt{socket}).
\end{itemize}

Mỗi chuỗi bao gồm ID syscall, dấu thời gian và siêu dữ liệu ngữ cảnh (ví dụ: ID tiến trình,
namespace). Với sự mất cân bằng lớp vốn có, các kỹ thuật lấy mẫu được áp dụng để ngăn ngừa
thiên lệch mô hình.

\subsubsection{Khai Thác Mẫu PrefixSpan: Trích Xuất Các Template Syscall Đáng Tin Cậy}
\label{subsubsec:prefixspan}

Để giảm thiểu dương tính giả và cải thiện hiệu quả phát hiện, DeSFAM tích hợp
PrefixSpan~\cite{prefixspan_pei}---một thuật toán khai thác mẫu tuần tự---trên tập con lành mạnh
của bộ dữ liệu DongTing. Điều này xác định các chuỗi con syscall thường xuyên đại diện cho hành
vi container hợp lệ.

Cơ Sở Dữ Liệu Mẫu Lành Mạnh kết quả được sử dụng để:
\begin{itemize}
  \item Khớp các cửa sổ syscall đến với các mẫu đã biết là tốt để phân loại lành mạnh sớm.
  \item Bỏ qua chấm điểm bất thường sâu cho các chuỗi khớp, giảm chi phí tính toán.
  \item Cung cấp khả năng truy xuất ngữ nghĩa cho các quyết định bảo mật.
\end{itemize}

\subsubsection{Kỹ Thuật Đặc Trưng: Cấu Trúc Hóa Chuỗi Syscall cho Phát Hiện}
\label{subsubsec:features}

SyscallAD áp dụng kỹ thuật đặc trưng đa chiều để chuyển đổi các luồng syscall thô thành các đầu
vào mô hình có thể hành động:

\begin{itemize}
  \item \textbf{Cửa Sổ Trượt:} Các syscall được phân đoạn thành các cửa sổ chồng lấp (độ
        dài~15, bước~3), nắm bắt cả các bất thường ngắn hạn và các mẫu hành vi mở rộng.
  \item \textbf{Tần Suất Danh Mục:} Các syscall được ánh xạ tới các danh mục chức năng (ví dụ:
        I/O tệp, bộ nhớ, mạng) và tần suất của chúng mỗi cửa sổ tạo thành chữ ký hành vi.
  \item \textbf{Đặc Trưng Thời Gian:} Các khoảng thời gian giữa syscall (delta thời gian) được
        tóm tắt (trung bình, độ lệch chuẩn, giá trị lớn nhất) để phát hiện các hành vi thời
        gian bất thường.
  \item \textbf{Khớp Mẫu Danh Sách Truy Cập:} Mỗi cửa sổ được kiểm tra chéo với các mẫu lành
        mạnh từ PrefixSpan; các kết khớp làm tăng độ tin cậy trong phân loại lành mạnh, giảm
        dương tính giả.
  \item \textbf{Tăng Cường Dữ Liệu:} Các chuỗi độc hại được sửa đổi tổng hợp (ví dụ: chèn
        hoặc hoán vị các syscall rủi ro cao như \texttt{clone}, \texttt{setuid}) để mô phỏng
        các kịch bản zero-day, cải thiện độ mạnh mẽ của mô hình.
\end{itemize}

\subsubsection{SyscallAD: Engine Học Sâu Kết Hợp và Phát Hiện Bất Thường}
\label{subsubsec:syscallad}

SyscallAD là một module phát hiện bất thường kết hợp được thiết kế để xác định các độ lệch hành
vi trong chuỗi syscall với độ chính xác và khả năng thích ứng cao.

\paragraph{Bộ Mã Hóa Biến Phân Được Tối Ưu Hóa (VAE)}
VAE sử dụng kiến trúc encoder-decoder ($32 \to 8 \to 32$ neurons) với các kích hoạt SELU,
dropout ($p = 0.2$), và chính quy hóa L2. Nó học cách tái tạo các cửa sổ syscall lành mạnh bằng
cách tối thiểu hóa Lỗi Tái Tạo Bình Phương Trung Bình (MSRE). Các bất thường được gắn cờ khi
MSRE vượt quá phân vị thứ 99.5 của các lỗi huấn luyện lành mạnh.

\paragraph{Isolation Forest (iForest)}
iForest cô lập các bất thường thông qua phân vùng đệ quy. Các chuỗi syscall hiếm gặp tạo ra các
cây cô lập nông hơn và điểm bất thường cao hơn. Tham số nhiễm là 2\% kiểm soát độ nhạy, cho
phép phát hiện các hành vi hiếm về mặt thống kê bị VAE bỏ qua.

\paragraph{Chiến Lược Hợp Nhất Tập Hợp}
SyscallAD kết hợp các điểm mạnh của VAE và Isolation Forest thông qua cơ chế hợp nhất điểm có
trọng số. Mỗi cửa sổ syscall $W_i$ được đánh giá bởi cả hai mô hình:
\begin{itemize}
  \item $A_{\text{VAE}}(W_i)$: Điểm bất thường dựa trên Lỗi Tái Tạo Bình Phương Trung Bình.
  \item $A_{\text{iForest}}(W_i)$: Điểm xuất phát từ độ sâu cô lập trong Isolation Forest.
\end{itemize}

Điểm cuối cùng được tính như sau:
\begin{equation}
  A(W_i) = \alpha \cdot A_{\text{VAE}}(W_i) + (1 - \alpha) \cdot A_{\text{iForest}}(W_i)
  \label{eq:ensemble}
\end{equation}
trong đó $\alpha = 0.7$ ưu tiên đóng góp của VAE. Giá trị này được xác nhận qua phân tích độ
nhạy trên $\alpha \in \{0.0, 0.25, 0.5, 0.7, 0.85, 1.0\}$ trên bộ dữ liệu DongTing, như thể
hiện trong Hình~\ref{fig:sensitivity}.

\begin{figure*}[tp]
  \centering
  \includegraphics[width=\linewidth]{figure6}
  \caption{Phân tích độ nhạy của trọng số hợp nhất $\alpha$ trên F1-Score và AUC.
           Hiệu suất đỉnh được quan sát tại $\alpha = 0.7$.}
  \label{fig:sensitivity}
\end{figure*}

Hình~\ref{fig:syscallad_output} trình bày nhật ký bất thường từ các container \texttt{test-mysql}
và \texttt{test-redis}. Các syscall lành mạnh như \texttt{futex}, \texttt{epoll\_wait} và
\texttt{wait4} được phân loại chính xác, trong khi các bất thường như \texttt{clone} và
\texttt{ioctl}, chỉ ra các nỗ lực thoát hoặc leo thang đặc quyền, được gắn cờ hiệu quả.

\begin{figure*}[tp]
  \centering
  \includegraphics[width=\linewidth]{figure5}
  \caption{Đầu ra SyscallAD từ DeSFAM Giai Đoạn~2: hành vi bình thường trong
           \texttt{test-mysql} tương phản với các phát hiện syscall bất thường trong
           \texttt{test-redis} (ví dụ: \texttt{clone}, \texttt{ioctl}).}
  \label{fig:syscallad_output}
\end{figure*}

\paragraph{Ngưỡng Động: Điều Chỉnh Ranh Giới Bất Thường Thích Ứng}
Để phù hợp với hành vi syscall đang thay đổi, SyscallAD sử dụng ngưỡng động dựa trên Trung Bình
Động Hàm Mũ (EMA), thay vì một điểm cắt cố định. Ngưỡng ban đầu $T_0$ được đặt ở phân vị thứ
99.5 của các lỗi tái tạo VAE trên dữ liệu xác nhận lành mạnh. Trong thời gian chạy, ngưỡng
được cập nhật sau mỗi cửa sổ syscall $W_i$ như sau:
\begin{equation}
  T_{t+1} = \max\bigl(T_{\min},\; \beta T_t + (1 - \beta) A(W_i)\bigr)
  \label{eq:threshold}
\end{equation}
trong đó $\beta = 0.9$ kiểm soát độ nhạy cập nhật và $T_{\min} = 0.5$ thực thi giới hạn dưới.

Phát hiện cuối cùng được định nghĩa là:
\begin{equation}
  \mathrm{Flag}(W_i) =
  \begin{cases}
    1, & \text{nếu } A(W_i) > T \\
    0, & \text{ngược lại}
  \end{cases}
  \label{eq:flag}
\end{equation}

\subsubsection{Giám Sát Thời Gian Thực với eBPF: Thu Thập Syscall Nhận Biết Container Chi Phí Thấp}
\label{subsubsec:ebpf_monitoring}

DeSFAM sử dụng các eBPF tracepoints để cho phép giám sát syscall chi tiết, chi phí thấp trực
tiếp ở mức kernel Linux. Chương trình BPF móc vào
\texttt{(raw\_syscalls:sys\_enter)} và \texttt{(raw\_syscalls:sys\_exit)}, cho phép chặn các lời
gọi syscall tại các điểm vào và thoát của chúng.

Để cô lập hoạt động syscall từ các khối lượng công việc container, DeSFAM ánh xạ động các PID
container bằng cách truy vấn Docker runtime và chèn chúng vào bảng băm BPF \texttt{pid\_map}.
Chỉ các syscall có nguồn gốc từ các PID này mới được ghi lại, đảm bảo theo dõi dành riêng cho
container mà không có nhiễu hệ thống không cần thiết.

Các sự kiện syscall được nắm bắt được làm giàu với siêu dữ liệu ngữ cảnh bao gồm: ID syscall,
dấu thời gian (độ phân giải nano giây), PID tiến trình và cha, tên lệnh, tên container và giá
trị trả về syscall. Chúng được gửi từ kernel đến không gian người dùng qua \texttt{ring\_buffer}
và được ghi vào các tệp theo container ở định dạng JSON. Đồng thời, các ID syscall được thêm vào
một deque cuộn (cửa sổ trượt kích thước~15) trong bộ nhớ cho mỗi container.
Hình~\ref{fig:ebpf_log} cho thấy một nhật ký syscall thời gian thực đại diện cho một container MySQL.

\begin{figure*}[tp]
  \centering
  \includegraphics[width=\linewidth]{figure7}
  \caption{Nhật ký syscall eBPF thời gian thực cho \texttt{kubepro-mysql-1}, hiển thị các sự
           kiện syscall có dấu thời gian với PID, tên container, người dùng và giá trị trả về.}
  \label{fig:ebpf_log}
\end{figure*}

\subsubsection{Tích Hợp DevOps: Tinh Chỉnh Mô Hình Liên Tục}
\label{subsubsec:devops}

Để giảm thiểu dương tính giả và thích nghi với các khối lượng công việc đang thay đổi, DeSFAM
tích hợp vòng phản hồi theo hướng DevOps trong Giai Đoạn~2. Các chuỗi syscall bất thường, cùng
với siêu dữ liệu (ví dụ: tên container, PID, dấu thời gian), được ghi nhật ký và trực quan hóa
để nhà phân tích xem xét. Các nhóm bảo mật có thể chú thích các phát hiện là dương tính thật
hoặc dương tính giả, với các dương tính giả đã xác nhận được thêm vào tập huấn luyện lành mạnh.

SyscallAD định kỳ huấn luyện lại tập hợp VAE + iForest của mình trên bộ dữ liệu tăng cường này,
cho phép nó học các biến thể lành mạnh mới được giới thiệu bởi các khối lượng công việc đặc thù
môi trường hoặc các trường hợp ngoại lệ.

Quy trình phát hiện đầy đủ được hình thức hóa trong Thuật toán~\ref{alg:phase2}.

\begin{algorithm}[htbp]
\caption{Phát Hiện Bất Thường Syscall Thời Gian Thực (SyscallAD)}
\label{alg:phase2}
\begin{algorithmic}[1]
  \Require eBPF hook, các mô hình VAE và iForest đã huấn luyện,
           các lỗi tái tạo xác nhận $E_{\text{val}}$
  \Ensure Cờ bất thường cho mỗi cửa sổ syscall
  \State Khởi tạo các cửa sổ trượt $W_i \leftarrow \emptyset$ cho $i \in \{5, 10, 15\}$
  \State Tải các mô hình VAE và iForest đã huấn luyện
  \State Tính $T_0 \leftarrow \text{percentile}(E_{\text{val}}, 99.5)$
  \State Đặt ngưỡng $T \leftarrow T_0$
  \While{luồng syscall đang hoạt động}
    \State Thu thập syscall $s_t$ tại thời điểm $t$
    \State Thêm $s_t$ vào mỗi $W_i$
    \For{mỗi cửa sổ đầy $W_i$}
      \State $f_i \leftarrow \text{extract\_features}(W_i)$
      \State $A_{\text{VAE}} \leftarrow \text{VAE\_score}(f_i)$
      \State $A_{\text{iForest}} \leftarrow \text{iForest\_score}(f_i)$
      \State $A(W_i) \leftarrow 0.7 \cdot A_{\text{VAE}} + 0.3 \cdot A_{\text{iForest}}$
      \If{$A(W_i) > T$}
        \State Gắn cờ $W_i$ là bất thường; leo thang lên Giai Đoạn~3
      \Else
        \State $T \leftarrow \max(T_{\min},\; \beta \cdot T + (1 - \beta) \cdot A(W_i))$
      \EndIf
      \State Trượt $W_i$ về phía trước
    \EndFor
  \EndWhile
\end{algorithmic}
\end{algorithm}

\subsection{Giai Đoạn 3: Thực Thi Syscall Thích Ứng và Quản Lý Chính Sách}
\label{subsec:phase3}

Thực Thi Syscall Thích Ứng là giai đoạn thứ ba của framework DeSFAM. Nó tập trung vào việc
thực thi động các chính sách syscall trong các môi trường container theo thời gian thực. Giai
đoạn này tích hợp chấm điểm rủi ro, phân tích ngữ cảnh và thông tin tình báo mối đe dọa bên
ngoài để cung cấp các phản hồi bảo mật thích ứng và chính xác.

\subsubsection{Cơ Chế Chấm Điểm Rủi Ro}
\label{subsubsec:risk_scoring}

DeSFAM định lượng mức độ đe dọa của mỗi chuỗi syscall bất thường $W_i$ bằng điểm rủi ro tổng
hợp được suy ra từ ba thành phần.

\paragraph{Ánh Xạ MITRE ATT\&CK}
Mỗi syscall trong $W_i$ được ánh xạ tới các kỹ thuật đối thủ liên quan (ví dụ: T1055 cho
\texttt{ptrace}) dựa trên bảng tra cứu được xác định trước (Bảng~\ref{tab:mitre_mapping}). Ánh
xạ ngữ cảnh này cho phép suy luận vượt ra ngoài các lời gọi riêng lẻ để nhận ra các mẫu tấn
công cấp cao.

\begin{table*}[tp]
  \centering
  \caption{Ví dụ Ánh Xạ Syscall-sang-Kỹ Thuật MITRE ATT\&CK.}
  \label{tab:mitre_mapping}
  \includegraphics[width=\linewidth]{table2}
\end{table*}

\paragraph{Mức Nghiêm Trọng Tích Lũy}
Điểm mức nghiêm trọng tích lũy tổng hợp các điểm mức nghiêm trọng riêng lẻ của các syscall
trong chuỗi:
\begin{equation}
  \text{Cumulative Severity} = \sum_{t \in W_i} \text{Severity}(t)
  \label{eq:severity}
\end{equation}

\paragraph{Khớp Mẫu Khai Thác}
Mỗi chuỗi $W_i$ được đánh giá đối chiếu với các chữ ký khai thác CVE đã biết. Nếu tìm thấy
kết khớp, một hình phạt rủi ro bổ sung được áp dụng:
\begin{equation}
  I(\text{Exploit Match}) =
  \begin{cases}
    1, & \text{nếu khớp} \\
    0, & \text{ngược lại}
  \end{cases}
  \label{eq:exploit_match}
\end{equation}

\paragraph{Điểm Rủi Ro Cuối Cùng}
Điểm rủi ro cuối cùng tích hợp cả mức nghiêm trọng và thông tin tình báo khai thác:
\begin{equation}
  \text{Risk Score} = \sum_{t \in W_i} \text{Severity}(t) + 5 \cdot I(\text{Exploit Match})
  \label{eq:risk_score}
\end{equation}

\paragraph{Logic Quyết Định Chính Sách}
Dựa trên điểm rủi ro được tính toán, DeSFAM áp dụng chiến lược thực thi ba cấp:
\begin{equation}
  \text{Action} =
  \begin{cases}
    \text{Block},       & \text{nếu Risk Score} \geq 8 \\
    \text{Rate-limit},  & \text{nếu } 5 \leq \text{Risk Score} < 8 \\
    \text{Log},         & \text{nếu Risk Score} < 5
  \end{cases}
  \label{eq:action}
\end{equation}

Các ngưỡng này được điều chỉnh theo kinh nghiệm và vẫn có thể cấu hình cho các tư thế rủi ro
đặc thù triển khai. Bảng~\ref{tab:risk_examples} trình bày các chuỗi mẫu và các hành động tương
ứng.

\begin{table}[htbp]
  \centering
  \caption{Điểm Rủi Ro Ví Dụ cho Các Chuỗi Lành Mạnh và Độc Hại.}
  \label{tab:risk_examples}
  \includegraphics[width=\linewidth]{table3}
\end{table}

\subsubsection{Thực Thi Dựa trên eBPF}
\label{subsubsec:ebpf_enforcement}

DeSFAM thực thi các chính sách syscall thông qua các chương trình eBPF được gắn vào LSM hooks.
Cụ thể, các chương trình này được tải và gắn bằng \texttt{bpf\_program\_attach\_lsm()} để chặn
các system call tại LSM hook \texttt{security\_syscall}.

\paragraph{Gắn LSM Hook}
Các chương trình eBPF được gắn động vào LSM hook \texttt{security\_syscall}:

\begin{lstlisting}[float=htbp, language=C, caption={Gắn LSM hook eBPF.}]
// eBPF Hook Attachment
struct bpf_program *prog =
    bpf_object__find_program_by_name(obj, "enforce_policy");
struct bpf_link *link = bpf_program_attach_lsm(prog);
\end{lstlisting}

Chữ ký hàm hook:
\begin{lstlisting}[float=htbp, language=C, caption={Chữ ký hàm LSM hook.}]
SEC("lsm/syscall")
int BPF_PROG(enforce_policy, int syscall_num,
             struct pt_regs *regs) {
    // Enforcement logic here
}
\end{lstlisting}

\paragraph{Cấu Trúc Policy Map}
Các hành động syscall được lưu trữ trong eBPF map dựa trên băm được lập chỉ mục bởi các định
danh container:

\begin{lstlisting}[float=htbp, language=C, caption={eBPF policy map theo khóa container.}]
struct container_key {
    u32 mntns;  // Mount namespace ID
    u32 pidns;  // PID namespace ID
};

struct bpf_map_def SEC("maps") policy_map = {
    .type       = BPF_MAP_TYPE_HASH,
    .key_size   = sizeof(struct container_key),
    .value_size = sizeof(u64),
    .max_entries = 144,
};
\end{lstlisting}

Map được ghim tại \texttt{/sys/fs/bpf/blocked\_syscalls} để cập nhật thời gian chạy qua
\texttt{policy\_loader}.

\paragraph{Cơ Chế Chặn Syscall}
Việc chặn thực sự các syscall xảy ra trong eBPF hook được gắn vào LSM
\texttt{security\_syscall}. Nguyên tắc thực thi quan trọng sử dụng ghi đè hàm kernel:
\begin{lstlisting}[float=htbp, language=C, caption={Chặn syscall qua ghi đè hàm kernel.}]
// Immediate syscall denial
bpf_override_return(regs, -EPERM);
\end{lstlisting}

Quy trình thực thi:
\begin{enumerate}
  \item Lấy ngữ cảnh container: \texttt{mntns = bpf\_get\_ns\_id(NS\_MNT);
        pidns = bpf\_get\_ns\_id(NS\_PID);}
  \item Truy vấn policy map: \texttt{u64 *action = bpf\_map\_lookup\_elem(\&policy\_map, \&key);}
  \item Thực thi hành động:
        \textbf{BLOCK}: Từ chối qua \texttt{bpf\_override\_return(regs, -EPERM)};
        \textbf{RATE\_LIMIT}: Cho phép $N$ lời gọi/phút qua bộ đếm nguyên tử;
        \textbf{LOG}: Ghi siêu dữ liệu qua ring buffer.
\end{enumerate}

\paragraph{Tự Phục Hồi Danh Sách Truy Cập qua Tái Tạo Seccomp}
Để đảm bảo sự tiến hóa chính sách lâu dài, DeSFAM kết hợp cơ chế tự phục hồi tự động tinh
chỉnh các danh sách truy cập syscall tĩnh dựa trên các lần chặn thời gian chạy lặp lại.

Các syscall bị chặn được theo dõi trong kernel bằng cách sử dụng LRU cache map gọi là
\texttt{probation\_map}:

\begin{lstlisting}[float=htbp, language=C, caption={Probation map để theo dõi các lần chặn lặp lại.}]
struct probation_key {
    struct container_key ckey;
    u32 syscall;  // Syscall number
};

SEC("maps") struct bpf_map_def probation_map = {
    .type       = BPF_MAP_TYPE_LRU_HASH,
    .key_size   = sizeof(struct probation_key),
    .value_size = sizeof(u32),  // Block counter
    .max_entries = 512,
};
\end{lstlisting}

Nếu một syscall bị chặn ba lần trở lên trong một cửa sổ có thể cấu hình và chưa được chú thích
là lành mạnh, kernel kích hoạt một daemon không gian người dùng qua
\texttt{bpf\_send\_signal(SIGUSR1)}. Daemon này giải quyết tên syscall, định vị và cập nhật hồ
sơ seccomp tương ứng của container, và loại bỏ syscall được gắn cờ khỏi danh sách truy cập. Vòng
phòng thủ thích ứng này đảm bảo rằng các container mới được triển khai kế thừa các chính sách
syscall được tăng cường, nhận biết khai thác mà không cần khởi động lại hay xem xét thủ công.

\paragraph{Mô Hình Chính Sách Hình Thức}
Cho $S$ là tập các syscall hợp lệ, và $M$, $N$ đại diện cho các namespace mount và PID. Chính
sách thực thi được định nghĩa là:
\begin{equation}
  P : S \times M \times N \to \{0, 1\}
  \label{eq:policy_func}
\end{equation}
\begin{equation}
  \text{Decision}(s) =
  \begin{cases}
    \text{Allow}, & \text{nếu } P(s, \mathit{mntns}, \mathit{pidns}) = 1 \\
    \text{Block},  & \text{ngược lại}
  \end{cases}
  \label{eq:decision}
\end{equation}

Logic nhận biết container này đảm bảo cô lập syscall và thực thi truy cập đặc quyền tối thiểu
theo namespace. Việc kết hợp sử dụng eBPF tracepoints và LSM hooks cho phép thực thi nhanh, phân
lớp với chi phí tối thiểu.

Quy trình thực thi trong Giai Đoạn~3 được hình thức hóa qua Thuật toán~\ref{alg:phase3}, và
quy trình thực thi đầy đủ được minh họa trong Hình~\ref{fig:phase3}.

\begin{figure}[htbp]
  \centering
  \includegraphics[width=\linewidth]{figure8}
  \caption{DeSFAM Giai Đoạn~3 -- Quy trình thực thi thích ứng. Các chuỗi syscall bất thường từ
           Giai Đoạn~2 bị chặn và xử lý thông qua eBPF và LSM hooks.}
  \label{fig:phase3}
\end{figure}

\begin{algorithm}[htbp]
\caption{Thực Thi Syscall Thích Ứng với Tự Tăng Cường}
\label{alg:phase3}
\begin{algorithmic}[1]
  \Require Chuỗi bất thường $W_i$, ánh xạ MITRE, cơ sở dữ liệu khai thác,
           ngữ cảnh container (\textit{mntns}, \textit{pidns}), eBPF map \texttt{policy\_map}
  \Ensure Hành động thực thi: ``Block'', ``Rate-limit'' hoặc ``Log''
  \State Khởi tạo $\mathit{Risk\_Score} \leftarrow 0$
  \For{mỗi syscall $t$ trong $W_i$}
    \State $\mathit{Risk\_Score} \leftarrow \mathit{Risk\_Score} + \text{Severity}(t)$
  \EndFor
  \If{$W_i$ khớp với một khai thác trong cơ sở dữ liệu}
    \State $\mathit{Risk\_Score} \leftarrow \mathit{Risk\_Score} + 5$
  \EndIf
  \If{$\mathit{Risk\_Score} \geq 8$}
    \State $\mathit{action} \leftarrow$ ``Block''
  \ElsIf{$5 \leq \mathit{Risk\_Score} < 8$}
    \State $\mathit{action} \leftarrow$ ``Rate-limit''
  \Else
    \State $\mathit{action} \leftarrow$ ``Log''
  \EndIf
  \For{mỗi syscall $t$ trong $W_i$}
    \State Cập nhật $\texttt{policy\_map}[(t, \mathit{mntns}, \mathit{pidns})] \leftarrow \mathit{action}$
    \State Thực thi $\mathit{action}$ trên $t$ qua eBPF
    \State Ghi nhật ký sự kiện $(t, \mathit{action}, \mathit{mntns}, \mathit{pidns})$
  \EndFor
  \If{$\mathit{action} = $ ``Block''}
    \For{mỗi syscall $t$ trong $W_i$}
      \State $\mathit{probation\_count}[t] \leftarrow \mathit{probation\_count}[t] + 1$
      \If{$\mathit{probation\_count}[t] \geq 3$ \textbf{and} $t$ chưa được chú thích là lành mạnh}
        \State Kích hoạt cập nhật danh sách truy cập seccomp cho syscall $t$
      \EndIf
    \EndFor
  \EndIf
  \State \Return $\mathit{action}$
\end{algorithmic}
\end{algorithm}

\subsubsection{Trường Hợp Sử Dụng: Nỗ Lực Leo Thang Đặc Quyền và Phòng Thủ CVE Thực Tế}
\label{subsubsec:usecase}

Xem xét một container chạy dưới chính sách hạn chế của DeSFAM được thiết kế để chặn các nỗ lực
leo thang đặc quyền. Một chuỗi tấn công điển hình---\texttt{clone()}, \texttt{setuid()},
\texttt{execve()}---bị chặn ở mức kernel.

Giả sử tiến trình hoạt động trong mount namespace $\mathit{mntns} = 1001$ và PID namespace
$\mathit{pidns} = 2001$. Khi gọi \texttt{setuid()}, DeSFAM tham khảo chính sách theo container
của nó: \texttt{policy\_map[1001, 2001] = [..., 1, 0, 1, ...]}, trong đó mỗi mục biểu thị hành
động cho một ID syscall ($0 = \text{BLOCK}$, $1 = \text{ALLOW}$). Vì \texttt{setuid()} được ánh
xạ tới 0, syscall bị từ chối. LSM hook eBPF, \texttt{realtime\_blocking}, thực thi quyết định
này trong kernel, trả về \texttt{-EPERM} và ghi nhật ký sự kiện để kiểm toán. Điều này đảm bảo
giảm thiểu ngay lập tức, trong kernel các nỗ lực khai thác mà không cần khởi động lại container
hay can thiệp không gian người dùng.
